%==============================================================================
% tento soubor pouzijte jako zaklad
% this file should be used as a base for the thesis
% (c) 2008 Michal Bidlo
% E-mail: bidlom AT fit vutbr cz
% Šablonu upravil / template edited by: Ing. Jaroslav Dytrych, dytrych@fit.vutbr.cz
%==============================================================================
% kodovaní: UTF-8 (zmena prikazem iconv, recode nebo cstocs)
% encoding: UTF-8 (you can change it by command iconv, recode or cstocs)
%------------------------------------------------------------------------------
% zpracování / processing: make, make pdf, make clean
%==============================================================================
% Soubory, které je nutné upravit: / Files which have to be edited:
%   xpolas34-ECS-20-literatura-bibliography.bib - literatura / bibliography
%   xpolas34-ECS-01-kapitoly-chapters.tex - obsah práce / the thesis content
%   xpolas34-ECS-30-prilohy-appendices.tex - přílohy / appendices
%==============================================================================
\documentclass[]{fitthesis} % bez zadání - pro začátek práce, aby nebyl problém s překladem
%\documentclass[english]{fitthesis} % without assignment - for the work start to avoid compilation problem
%\documentclass[zadani]{fitthesis} % odevzdani do wisu - odkazy jsou barevné
%\documentclass[english,zadani]{fitthesis} % for submission to the IS FIT - links are color
%\documentclass[zadani,print]{fitthesis} % pro tisk - odkazy jsou černé
%\documentclass[english,zadani,print]{fitthesis} % for the print - links are black
% * Je-li prace psana v anglickem jazyce, je zapotrebi u tridy pouzit 
%   parametr english nasledovne:
%   If thesis is written in english, it is necessary to use 
%   parameter english as follows:
%      \documentclass[english]{fitthesis}
% * Je-li prace psana ve slovenskem jazyce, je zapotrebi u tridy pouzit 
%   parametr slovak nasledovne:
%      \documentclass[slovak]{fitthesis}

% Základní balíčky jsou dole v souboru šablony fitthesis.cls
% Basic packages are at the bottom of template file fitthesis.cls
%zde muzeme vlozit vlastni balicky / you can place own packages here

%---rm---------------
\renewcommand{\rmdefault}{lmr}%zavede Latin Modern Roman jako rm / set Latin Modern Roman as rm
%---sf---------------
\renewcommand{\sfdefault}{qhv}%zavede TeX Gyre Heros jako sf
%---tt------------
\renewcommand{\ttdefault}{lmtt}% zavede Latin Modern tt jako tt

% vypne funkci šablony, která automaticky nahrazuje uvozovky,
% aby nebyly prováděny nevhodné náhrady v popisech API apod.
% disables function of the template which replaces quotation marks
% to avoid unnecessary replacements in the API descriptions etc.
\csdoublequotesoff

% =======================================================================
% balíček "hyperref" vytváří klikací odkazy v pdf, pokud tedy použijeme pdflatex
% problém je, že balíček hyperref musí být uveden jako poslední, takže nemůže
% být v šabloně
% "hyperref" package create clickable links in pdf if you are using pdflatex.
% Problem is that this package have to be introduced as the last one so it 
% can not be placed in the template file.
\ifWis
\ifx\pdfoutput\undefined % nejedeme pod pdflatexem / we are not using pdflatex
\else
  \usepackage{color}
  \usepackage[unicode,colorlinks,hyperindex,plainpages=false,pdftex]{hyperref}
  \definecolor{links}{rgb}{0.4,0.5,0}
  \definecolor{anchors}{rgb}{1,0,0}
  \def\AnchorColor{anchors}
  \def\LinkColor{links}
  \def\pdfBorderAttrs{/Border [0 0 0] }  % bez okrajů kolem odkazů / without margins around links
  \pdfcompresslevel=9
\fi
\else % pro tisk budou odkazy, na které se dá klikat, černé / for the print clickable links will be black
\ifx\pdfoutput\undefined % nejedeme pod pdflatexem / we are not using pdflatex
\else
  \usepackage{color}
  \usepackage[unicode,colorlinks,hyperindex,plainpages=false,pdftex,urlcolor=black,linkcolor=black,citecolor=black]{hyperref}
  \definecolor{links}{rgb}{0,0,0}
  \definecolor{anchors}{rgb}{0,0,0}
  \def\AnchorColor{anchors}
  \def\LinkColor{links}
  \def\pdfBorderAttrs{/Border [0 0 0] } % bez okrajů kolem odkazů / without margins around links
  \pdfcompresslevel=9
\fi
\fi
% Řešení problému, kdy klikací odkazy na obrázky vedou za obrázek
% This solves the problems with links which leads after the picture
\usepackage[all]{hypcap}

% Informace o práci/projektu / Information about the thesis
%---------------------------------------------------------------------------
\projectinfo{
  %Prace / Thesis
  project=BP,            %typ prace BP/SP/DP/DR  / thesis type (SP = term project)
  year=2017,             %rok odevzdání / year of submission
  date=\today,           %datum odevzdani / submission date
  %Nazev prace / thesis title
  title.cs={Komponentní systém pro herní grafický engine},  %nazev prace v cestine ci slovenstine (dle zadani) / thesis title in czech language (according to assignment)
  title.en={Graphic Engine Based on Entity Component System}, %nazev prace v anglictine / thesis title in english
  %Autor / Author
  author={Tomáš Polášek},   %cele jmeno a prijmeni autora / full name and surname of the author
  author.name={Tomáš},   %jmeno autora (pro citaci) / author name (for reference) 
  author.surname={Polášek},   %prijmeni autora (pro citaci) / author surname (for reference) 
  %author.title.p=Bc., %titul pred jmenem (nepovinne) / title before the name (optional)
  %author.title.a=PhD, %titul za jmenem (nepovinne) / title after the name (optional)
  %Ustav / Department
  department=UPGM, % doplnte prislusnou zkratku dle ustavu na zadani: UPSY/UIFS/UITS/UPGM
  %                  fill in appropriate abbreviation of the department according to assignment: UPSY/UIFS/UITS/UPGM
  %Skolitel / supervisor
  supervisor=Michal Španěl, %cele jmeno a prijmeni skolitele / full name and surname of the supervisor
  supervisor.name={Michal},   %jmeno skolitele (pro citaci) / supervisor name (for reference) 
  supervisor.surname={Španěl},   %prijmeni skolitele (pro citaci) / supervisor surname (for reference) 
  supervisor.title.p=,   %titul pred jmenem (nepovinne) / title before the name (optional)
  supervisor.title.a={Ing., Ph.D.},    %titul za jmenem (nepovinne) / title after the name (optional)
  %Klicova slova, abstrakty, prohlaseni a podekovani je mozne definovat 
  %bud pomoci nasledujicich parametru nebo pomoci vyhrazenych maker (viz dale)
  %Keywords, abstracts, declaration and acknowledgement can be defined by following 
  %parameters or using dedicated macros (see below)
  %===========================================================================
  %Klicova slova / keywords
  %keywords.cs={Klíčová slova v českém jazyce.}, %klicova slova v ceskem ci slovenskem jazyce
  %                                              keywords in czech or slovak language
  %keywords.en={Klíčová slova v anglickém jazyce.}, %klicova slova v anglickem jazyce / keywords in english
  %Abstract
  %abstract.cs={Výtah (abstrakt) práce v českém jazyce.}, % abstrakt v ceskem ci slovenskem jazyce
  %                                                         abstract in czech or slovak language
  %abstract.en={Výtah (abstrakt) práce v anglickém jazyce.}, % abstrakt v anglickem jazyce / abstract in english
  %Prohlaseni / Declaration
  %declaration={Prohlašuji, že jsem tuto bakalářskou práci vypracoval samostatně pod vedením pana ...},
  %Podekovani (nepovinne) / Acknowledgement (optional)
  %acknowledgment={Zde je možné uvést poděkování vedoucímu práce a těm, kteří poskytli odbornou pomoc.} % nepovinne
  %acknowledgment={Here it is possible to express thanks to the supervisor and to the people which provided professional help.} % optional
}

%Abstrakt (cesky, slovensky ci anglicky) / Abstract (in czech, slovak or english)
\abstract[cs]{Cílem této práce je návrh a implementace knihovny pro správu entit tvořených pomocí kompozice z menších částí. Součástí práce je analýza dnes používaných technik pro práci s entitami -- \emph{objektově orientovaná dědičnost}, \emph{objektově orientovaná kompozice} a \emph{datově orientovaná kompozice} -- a jejich porovnání. \\
Práce představuje návrh komponentního systému a zdůvodnění vybraných rozhodnutí. Dále je představena implementace v programovacím jazyce \emph{C++}. Nakonec je ukázáno možné využití vytvořené knihovny a návrh hry za použití \emph{datově orientované kompozice}. \\
Vytvořené řešení umožňuje dynamickou kompozici entit za běhu aplikace z komponent definovaných za jejího překladu. Navržený systém je možné začlenit do již existujících projektů, umožňuje paralelní přístup a je uživatelsky rozšiřitelný. \\
Výsledná knihovna umožňuje návrh aplikací za použití \emph{datově orientované kompozice}, bez nutnosti implementace vlastního systému.}
\abstract[en]{The goal of this bachelor thesis is to design and implement a library for management of entities created using \emph{dat-oriented composition}. Analysis of the ways entities are used in todays software is performed -- \emph{object-oriented inheritance}, \emph{object-oriented composition} and \emph{data-oriented composition}. \\
Possible design of such composition system is presented and chosen decisions are discussed. The aforementioned design is then implemented in the \emph{C++} programming language. In the end a demo game implementation is presented using the \emph{data-oriented composition} pattern. \\
Presented solution allows dynamic composition of entities during application runtime from components defined at compile time. System is easily incorporated into existing projects, allows parallel access and can be extended by the user. \\
Resulting library allows experimentation and designing  of applications using \emph{data-oriented composition} pattern, without the need to implement the whole system.}

%Klicova slova (cesky, slovensky ci anglicky) / Keywords (in czech, slovak or english)
\keywords[cs]{Komponentní systém, Entitní systém, ECS, Kompozice, Data-oriented design, Paralelizmus, C++}
\keywords[en]{Component system, Entity-component system, ECS, Composition, Data-oriented design, Parallel, C++}

%Prohlaseni (u anglicky psane prace anglicky, u slovensky psane prace slovensky)
%Declaration (for thesis in english should be in english)
\declaration{Prohlašuji, že jsem tuto bakalářskou práci vypracoval samostatně pod vedením pana X...
Další informace mi poskytli...
Uvedl jsem všechny literární prameny a publikace, ze kterých jsem čerpal.}

% \declaration{Hereby I declare that this bachelor's thesis was prepared as an original author’s work under the supervision of Mr. X
% The supplementary information was provided by Mr. Y
% All the relevant information sources, which were used during preparation of this thesis, are properly cited and included in the list of references.}

%Podekovani (nepovinne, nejlepe v jazyce prace) / Acknowledgement (optional, ideally in the language of the thesis)
\acknowledgment{V této sekci je možno uvést poděkování vedoucímu práce a těm, kteří poskytli odbornou pomoc
(externí zadavatel, konzultant, apod.).}
%\acknowledgment{Here it is possible to express thanks to the supervisor and to the people which provided professional help
%(external submitter, consultant, etc.).}

% řeší první/poslední řádek odstavce na předchozí/následující stránce
% solves first/last row of the paragraph on the previous/next page
\clubpenalty=10000
\widowpenalty=10000

\begin{document}
  % Vysazeni titulnich stran / Typesetting of the title pages
  % ----------------------------------------------
  \maketitle
  % Obsah
  % ----------------------------------------------
  \tableofcontents
  
  % Seznam obrazku a tabulek (pokud prace obsahuje velke mnozstvi obrazku, tak se to hodi)
  % List of figures and list of tables (if the thesis contains a lot of pictures, it is good)
\ifczech
  \renewcommand\listfigurename{Seznam obrázků}
\fi
\ifslovak
  \renewcommand\listfigurename{Zoznam obrázkov}
\fi

  % \listoffigures
\ifczech
  \renewcommand\listtablename{Seznam tabulek}
\fi
\ifslovak
  \renewcommand\listtablename{Zoznam tabuliek}
\fi

  % \listoftables 

  % vynechani stranky v oboustrannem rezimu
  % Skip the page in the two-sided mode
  \iftwoside
    \cleardoublepage
  \fi

  % Text prace / Thesis text
  % ----------------------------------------------
  %=========================================================================
% (c) Michal Bidlo, Bohuslav Křena, 2008


\chapter{Úvod}
\todo{Shrnutí cílů}
\blind[1]
\todo{Teorie, proč se tím zaobírat}
\blind[1]
\todo{Návrh a implementace}
\blind[1]
\todo{Vyhodnocení}
\blind[1]


\chapter{Teoretický rozbor}
\todo{Popis kapitoly}
\blind[1]

\section{Požadavky na návrh software}
\todo{Délka, složitost. Potřeba oddělení částí.}
\blind[2]

\section{Datově orientovaný návrh}
\todo{Proč se používá, architektura PC.}
\blind[4]

\section{Reprezentace entit a jejich chování}
\todo{Co to je entita, chování}
\blind[2]

\subsection{Objektově orientovaná hierarchie}
\todo{Popis, příklad, vyhodnocení výhod/nevýhod}
\begin{figure}[H]
	\tmpframe{\includegraphics[width=\linewidth]{TODO-image}}
	\caption{Příklad objektově orientované hierarchie. \todo{Vytvořit}}
\end{figure}
\begin{figure}[H]
	\tmpframe{\includegraphics[width=\linewidth]{TODO-image}}
	\caption{Obsah paměti u objektově orientovaných hierarchií. \todo{Vytvořit}}
\end{figure}
\blind[3]

\subsection{Objektově orientovaná kompozice}
\todo{Popis, příklad, vyhodnocení výhod/nevýhod}
\begin{figure}[H]
	\tmpframe{\includegraphics[width=\linewidth]{TODO-image}}
	\caption{Příklad objektově orientované kompozice. \todo{Vytvořit}}
\end{figure}
\begin{figure}[H]
	\tmpframe{\includegraphics[width=\linewidth]{TODO-image}}
	\caption{Obsah paměti u objektově orientované kompozice. \todo{Vytvořit}}
\end{figure}
\blind[3]

\subsection{Datově orientovaná kompozice}
\todo{Popis, příklad, vyhodnocení výhod/nevýhod}
\begin{figure}[H]
	\tmpframe{\includegraphics[width=\linewidth]{TODO-image}}
	\caption{Příklad datově orientované kompozice. \todo{Vytvořit}}
\end{figure}
\begin{figure}[H]
	\tmpframe{\includegraphics[width=\linewidth]{TODO-image}}
	\caption{Obsah paměti u datově orientované kompozice. \todo{Vytvořit}}
\end{figure}
\blind[4]

\section{Komponentní systémy}
\todo{Zaklad, co to je.}
\blind[1]

\subsection{Historie vývoje}
\blind[3]

\subsection{Koncepty}
\todo{Představení základní architektury návrhového vzoru - entity, komponenty, systémy}
\begin{figure}[H]
	\tmpframe{\includegraphics[width=\linewidth]{TODO-image}}
	\caption{Entita reprezentována jako řádek v databázi. \todo{Vytvořit}}
\end{figure}
\blind[5]

\subsection{Výhody zvolené reprezentace}
\todo{Lokalita dat, instrukcí. Cache. Dynamičnost, prototypování. Serializace. Embedded skriptovací jazyky.}
\blind[4]

\section{Paralelizmus}
\todo{Důležitost, hardware - cache, popis možností paralelizace.}
\blind[2]
\begin{figure}[H]
	\tmpframe{\includegraphics[width=\linewidth]{TODO-image}}
	\caption{Paralelizmus uvnitř systémů. \todo{Vytvořit}}
\end{figure}
\blind[1]
\begin{figure}[H]
	\tmpframe{\includegraphics[width=\linewidth]{TODO-image}}
	\caption{Paralelizmus na úrovni systémů. \todo{Vytvořit}}
\end{figure}
\blind[1]
%https://software.intel.com/en-us/articles/designing-the-framework-of-a-parallel-game-engine

\chapter{Návrh komponentního systému}
\todo{Postup seznámení s návrhem top -> down}
\begin{figure}[H]
	\tmpframe{\includegraphics[width=\linewidth]{TODO-image}}
	\caption{Návrhový diagram komponentního systému. \todo{Vytvořit}}
\end{figure}
\blind[1]

\section{Reprezentace entit}
\todo{Entita, id, generace. Alternativy a důvody.}
\blind[3]

\section{Komponenty a jejich nosiče}
\todo{Komponenta POD. Trivial constructor, copy. Paměť a typy nosičů. Typy komponent - tag, normal.}
\blind[4]

\section{Systémy a skupiny}
\todo{Volnost použití systémů. Uživatelem definované chování. Filtry - require, reject. Cache pomocí skupin entit.}
\blind[5]

\section{Paralelní přístup}
\todo{Podporované způsoby paralelizace, návrh, zdůvodnění. Entity paralellization. System parallelization. Diskuze na téma, zda je paralelní přístup potřeba.}
\blind[5]

\section{Přehled komunikace}
\todo{High-level design, propojení. Universe, action-cache.}
\begin{figure}[H]
	\tmpframe{\includegraphics[width=\linewidth]{TODO-image}}
	\caption{Přehled návrhu komponentního systému a komunikace mezi jeho částmi. \todo{Vytvořit - HLB diagram}}
\end{figure}
\blind[3]

\section{Tok řízení}
\todo{Popsat předpokládaný tok řízení, jak se bude používat. Registrace. Iterace. Refresh.}
\begin{figure}[H]
	\tmpframe{\includegraphics[width=\linewidth]{TODO-image}}
	\caption{Digram toku řízení komponentního systému. \todo{Vytvořit - Control flow diagram}}
\end{figure}
\blind[3]


\chapter{Implementace}
\todo{Shrnutí návrhu, zdůvodnění výběru implementačního jazyku C++}
\blind[2]

\section{Použité knihovny a přenositelnost}
\todo{Zdůvodnění nepoužití knihoven. C++14, Linux, Windows.}
\blind[2]

\section{Správa entit}
\todo{Potřebné informace, implementace, datové struktury, optimalizace}
\blind[3]

\section{Komponenty}
\todo{Rekapitulace vlastností, které komponenta musi splnit. Registrace, generování identifikátorů. Uživatelské nosiče}
\blind[3]

\section{Systémy a skupiny}
\todo{Super třída, filtry, přístup - iterace. Přidání systému, volitelná skupina - added/removed.}
\blind[3]

\section{Obnovovací fáze}
\todo{Pořadí refresh. Obnovení aktuality skupin.}
\blind[2]

\section{Podpora paralelizmu}
\todo{2 způsoby paralelizmu - dělení skupin, change-sety}
\blind[3]


\chapter{Vyhodnocení}
\todo{Shrnutí, popis co bude v této kapitole}
\blind[1]

\section{Testování knihovny}
\todo{Unit testing, benchmarky}
\blind[2]

\section{Testované sestavy}
\todo{Hardware sestavy - PC, notebook. Tabulka identifikace}
\blind[1]

\section{Výkonnostní testy}
\todo{Výkon při různých zatíženích - procenta přidaných/odebraných. Grafy.}
\blind[1]
\todo{Porovnání nosičů komponent. Grafy.}
\blind[1]
\todo{Seriové vs paralelní - počet vláken. Grafy.}
\blind[1]

\section{Porovnání}
\todo{Výkonostní porovnání oproti jiným ECS knihovnám, Unity? Grafy.}
\blind[2]

\section{Demo hra}
\todo{Důvod, typ hry}
\blind[1]

\subsection{Návrh}
\todo{Popis principů hry}
\blind[1]
\todo{Komponentní návrh}
\blind[1]

\subsection{Implementace}
\todo{Použité knihovny, základní popis implementace}
\blind[2]

\subsection{Výsledek}
\todo{Poznatky a výsledky získané ze hry. Grafy?}
\blind[2]


\chapter{Závěr}
\todo{Využití komponentních systémů a aktuální stav} 
\blind[1]
\todo{Návrh, prototypy, priority} 
\blind[1]
\todo{Implementace, rozšiřitelnost, přenositelnost}
\blind[1]
\todo{Experimenty, srovnání, demo hra}
\blind[1]
\todo{Možnosti rozšíření}
\blind[1]

%\begin{figure}
%	\tmpframe{\includegraphics[width=\linewidth]{TODO-image}}
%	\caption{Obrazek \todo{Obrazek}}
%\end{figure}

%=========================================================================
 % viz. obsah.tex / see obsah.tex

  % Pouzita literatura / Bibliography
  % ----------------------------------------------
\ifslovak
  \makeatletter
  \def\@openbib@code{\addcontentsline{toc}{chapter}{Literatúra}}
  \makeatother
  \bibliographystyle{bib-styles/czechiso}
\else
  \ifczech
    \makeatletter
    \def\@openbib@code{\addcontentsline{toc}{chapter}{Literatura}}
    \makeatother
    \bibliographystyle{bib-styles/czechiso}
  \else 
    \makeatletter
    \def\@openbib@code{\addcontentsline{toc}{chapter}{Bibliography}}
    \makeatother
    \bibliographystyle{bib-styles/englishiso}
  %  \bibliographystyle{alpha}
  \fi
\fi
  \begin{flushleft}
  \bibliography{xpolas34-ECS-20-literatura-bibliography}
  \end{flushleft}

  % vynechani stranky v oboustrannem rezimu
  % Skip the page in the two-sided mode
  \iftwoside
    \cleardoublepage
  \fi

  % Prilohy / Appendices
  % ---------------------------------------------
  \appendix
\ifczech
  \renewcommand{\appendixpagename}{Přílohy}
  \renewcommand{\appendixtocname}{Přílohy}
  \renewcommand{\appendixname}{Příloha}
\fi
\ifslovak
  \renewcommand{\appendixpagename}{Prílohy}
  \renewcommand{\appendixtocname}{Prílohy}
  \renewcommand{\appendixname}{Príloha}
\fi
  \appendixpage

% vynechani stranky v oboustrannem rezimu
% Skip the page in the two-sided mode
\iftwoside
  \cleardoublepage
\fi
  
\ifslovak
%  \section*{Zoznam príloh}
%  \addcontentsline{toc}{section}{Zoznam príloh}
\else
  \ifczech
%    \section*{Seznam příloh}
%    \addcontentsline{toc}{section}{Seznam příloh}
  \else
%    \section*{List of Appendices}
%    \addcontentsline{toc}{section}{List of Appendices}
  \fi
\fi
  \startcontents[chapters]
  % seznam příloh / list of appendices
  % \printcontents[chapters]{l}{0}{\setcounter{tocdepth}{2}}
  
  % vynechani stranky v oboustrannem rezimu
  \iftwoside
    \cleardoublepage
  \fi
  % Tento soubor nahraďte vlastním souborem s přílohami (nadpisy níže jsou pouze pro příklad)
% This file should be replaced with your file with an appendices (headings below are examples only)

% Umístění obsahu paměťového média do příloh je vhodné konzultovat s vedoucím
% Placing of table of contents of the memory media here should be consulted with a supervisor
%\chapter{Obsah přiloženého paměťového média}

%\chapter{Manuál}

%\chapter{Konfigurační soubor} % Configuration file

%\chapter{RelaxNG Schéma konfiguračního souboru} % Scheme of RelaxNG configuration file

%\chapter{Plakát} % poster

\chapter{Jak pracovat s touto šablonou}
\label{jak}

V této kapitole je uveden popis jednotlivých částí šablony, po kterém následuje stručný návod, jak s touto šablonou pracovat. 

Jedná se o přechodnou verzi šablony. Nová verze bude zveřejněna do konce roku 2016 a bude navíc obsahovat nové pokyny ke správnému využití šablony, závazné pokyny k~vypracování bakalářských a diplomových prací (rekapitulace pokynů, které jsou dostupné na~webu) a nezávazná doporučení od vybraných vedoucích. Jediné soubory, které se v nové verzi změní, budou projekt-01-kapitoly-chapters.tex a projekt-30-prilohy-appendices.tex, jejichž obsah každý student vymaže a nahradí vlastním. Šablonu lze tedy bez problémů využít i~v~současné verzi.

\section*{Popis částí šablony}

Po rozbalení šablony naleznete následující soubory a adresáře:
\begin{DESCRIPTION}
  \item [bib-styles] Styly literatury (viz níže). 
  \item [obrazky-figures] Adresář pro Vaše obrázky. Nyní obsahuje placeholder.pdf (tzv. TODO obrázek, který lze použít jako pomůcku při tvorbě technické zprávy), který se s prací neodevzdává. Název adresáře je vhodné zkrátit, aby byl jen ve zvoleném jazyce.
  \item [template-fig] Obrázky šablony (znak VUT).
  \item [fitthesis.cls] Šablona (definice vzhledu).
  \item [Makefile] Makefile pro překlad, počítání normostran, sbalení apod. (viz níže).
  \item [projekt-01-kapitoly-chapters.tex] Soubor pro Váš text (obsah nahraďte).
  \item [projekt-20-literatura-bibliography.bib] Seznam literatury (viz níže).
  \item [projekt-30-prilohy-appendices.tex] Soubor pro přílohy (obsah nahraďte).
  \item [projekt.tex] Hlavní soubor práce -- definice formálních částí.
\end{DESCRIPTION}

Výchozí styl literatury (czechiso) je od Ing. Martínka, přičemž anglická verze (englishiso) je jeho překladem s drobnými modifikacemi. Oproti normě jsou v něm určité odlišnosti, ale na FIT je dlouhodobě akceptován. Alternativně můžete využít styl od Ing. Radima Loskota nebo od Ing. Radka Pyšného\footnote{BP Ing. Radka Pyšného \url{http://www.fit.vutbr.cz/study/DP/BP.php?id=7848}}. Alternativní styly obsahují určitá vylepšení, ale zatím nebyly řádně otestovány větším množstvím uživatelů. Lze je považovat za beta verze pro zájemce, kteří svoji práci chtějí mít dokonalou do detailů a neváhají si nastudovat detaily správného formátování citací, aby si mohli ověřit, že je vysázený výsledek v pořádku.

Makefile kromě překladu do PDF nabízí i další funkce:
\begin{itemize}
  \item přejmenování souborů (viz níže),
  \item počítání normostran,
  \item spuštění vlny pro doplnění nezlomitelných mezer,
  \item sbalení výsledku pro odeslání vedoucímu ke kontrole (zkontrolujte, zda sbalí všechny Vámi přidané soubory, a případně doplňte).
\end{itemize}

Nezapomeňte, že vlna neřeší všechny nezlomitelné mezery. Vždy je třeba manuální kontrola, zda na konci řádku nezůstalo něco nevhodného -- viz Internetová jazyková příručka\footnote{Internetová jazyková příručka \url{http://prirucka.ujc.cas.cz/?id=880}}.

\paragraph {Pozor na číslování stránek!} Pokud má obsah 2 strany a na 2. jsou jen \uv{Přílohy} a~\uv{Seznam příloh} (ale žádná příloha tam není), z nějakého důvodu se posune číslování stránek o 1 (obsah \uv{nesedí}). Stejný efekt má, když je na 2. či 3. stránce obsahu jen \uv{Literatura} a~je možné, že tohoto problému lze dosáhnout i jinak. Řešení je několik (od~úpravy obsahu, přes nastavení počítadla až po sofistikovanější metody). \textbf{Před odevzdáním proto vždy překontrolujte číslování stran!}


\section*{Doporučený postup práce se šablonou}

\begin{enumerate}
  \item \textbf{Zkontrolujte, zda máte aktuální verzi šablony.} Máte-li šablonu z předchozího roku, na stránkách fakulty již může být novější verze šablony s~aktualizovanými informacemi, opravenými chybami apod.
  \item \textbf{Zvolte si jazyk}, ve kterém budete psát svoji technickou zprávu (česky, slovensky nebo anglicky) a svoji volbu konzultujte s vedoucím práce (nebyla-li dohodnuta předem). Pokud Vámi zvoleným jazykem technické zprávy není čeština, nastavte příslušný parametr šablony v souboru projekt.tex (např.: \verb|documentclass[english]{fitthesis}| a přeložte prohlášení a poděkování do~angličtiny či slovenštiny.
  \item \textbf{Přejmenujte soubory.} Po rozbalení je v šabloně soubor projekt.tex. Pokud jej přeložíte, vznikne PDF s technickou zprávou pojmenované projekt.pdf. Když vedoucímu více studentů pošle projekt.pdf ke kontrole, musí je pracně přejmenovávat. Proto je vždy vhodné tento soubor přejmenovat tak, aby obsahoval Váš login a (případně zkrácené) téma práce. Vyhněte se však použití mezer, diakritiky a speciálních znaků. Vhodný název tedy může být např.: \uv{xlogin00-Cisteni-a-extrakce-textu.tex}. K přejmenování můžete využít i přiložený Makefile:
\begin{verbatim}
make rename NAME=xlogin00-Cisteni-a-extrakce-textu
\end{verbatim}
  \item Vyplňte požadované položky v souboru, který byl původně pojmenován projekt.tex, tedy typ, rok (odevzdání), název práce, svoje jméno, ústav (dle zadání), tituly a~jméno vedoucího, abstrakt, klíčová slova a další formální náležitosti.
  \item Nahraďte obsah souborů s kapitolami práce, literaturou a přílohami obsahem svojí technické zprávy. Jednotlivé přílohy či kapitoly práce může být výhodné uložit do~samostatných souborů -- rozhodnete-li se pro toto řešení, je doporučeno zachovat konvenci pro názvy souborů, přičemž za číslem bude následovat název kapitoly. 
  \item Nepotřebujete-li přílohy, zakomentujte příslušnou část v projekt.tex a příslušný soubor vyprázdněte či smažte. Nesnažte se prosím vymyslet nějakou neúčelnou přílohu jen proto, aby daný soubor bylo čím naplnit. Vhodnou přílohou může být obsah přiloženého paměťového média.
  \item Nascanované zadání uložte do souboru zadani.pdf a povolte jeho vložení do práce parametrem šablony v projekt.tex (\verb|documentclass[zadani]{fitthesis}|).
  \item Nechcete-li odkazy tisknout barevně (tedy červený obsah -- bez konzultace s vedoucím nedoporučuji), budete pro tisk vytvářet druhé PDF s tím, že nastavíte parametr šablony pro tisk: (\verb|documentclass[zadani,print]{fitthesis}|).  Barevné logo se nesmí tisknout černobíle!
  \item Vzor desek, do kterých bude práce vyvázána, si vygenerujte v informačním systému fakulty u zadání. Pro disertační práci lze zapnout parametrem v šabloně (více naleznete v souboru fitthesis.cls).
  \item Nezapomeňte, že zdrojové soubory i (obě verze) PDF musíte odevzdat na CD či jiném médiu přiloženém k technické zprávě.
\end{enumerate}

\subsection*{Pokyny pro oboustranný tisk}
\begin{itemize}
\item Zapíná se parametrem šablony: \verb|\documentclass[twoside]{fitthesis}|
\item Po vytištění oboustranného listu zkontrolujte, zda je při prosvícení sazební obrazec na obou stranách na stejné pozici. Méně kvalitní tiskárny s duplexní jednotkou mají často posun o 1--3 mm. Toto může být u některých tiskáren řešitelné tak, že vytisknete nejprve liché stránky, pak je dáte do stejného zásobníku a vytisknete sudé.
\item Za titulním listem, obsahem, literaturou, úvodním listem příloh, seznamem příloh a případnými dalšími seznamy je třeba nechat volnou stránku, aby následující část začínala na liché stránce (\textbackslash cleardoublepage).
\item  Konečný výsledek je nutné pečlivě překontrolovat.
\end{itemize}


\subsection*{Užitečné nástroje}
\label{nastroje}

Následující seznam není výčtem všech využitelných nástrojů. Máte-li vyzkoušený osvědčený nástroj, neváhejte jej využít. Pokud však nevíte, který nástroj si zvolit, můžete zvážit některý z následujících:

\begin{description}
	\item[\href{http://miktex.org/download}{MikTeX}] \LaTeX{} pro Windows -- distribuce s jednoduchou instalací a vynikající automatizací stahování balíčků.
	\item[\href{http://texstudio.sourceforge.net/}{TeXstudio}] Přenositelné opensource GUI pro \LaTeX{}.  Ctrl+klik umožňuje přepínat mezi zdrojovým textem a PDF. Má integrovanou kontrolu pravopisu, zvýraznění syntaxe apod. Pro jeho využití je nejprve potřeba nainstalovat MikTeX.
	\item[\href{http://jabref.sourceforge.net/download.php}{JabRef}] Pěkný a jednoduchý program v Javě pro správu souborů s bibliografií (literaturou). Není potřeba se nic učit -- poskytuje jednoduché okno a formulář pro editaci položek.
	\item[\href{https://inkscape.org/en/download/}{InkScape}] Přenositelný opensource editor vektorové grafiky (SVG i PDF). Vynikající nástroj pro tvorbu obrázků do odborného textu. Jeho ovládnutí je obtížnější, ale výsledky stojí za to.
	\item[\href{https://git-scm.com/}{GIT}] Vynikající pro týmovou spolupráci na projektech, ale může výrazně pomoci i jednomu autorovi. Umožňuje jednoduché verzování, zálohování a přenášení mezi více počítači.
	\item[\href{http://www.overleaf.com/}{Overleaf}] Online nástroj pro \LaTeX{}. Přímo zobrazuje náhled a umožňuje jednoduchou spolupráci (vedoucí může průběžně sledovat psaní práce), vyhledávání ve zdrojovém textu kliknutím do PDF, kontrolu pravopisu apod. Zdarma jej však lze využít pouze s určitými omezeními (někomu stačí na disertaci, jiný na ně může narazit i při psaní bakalářské práce) a pro dlouhé texty je pomalejší.
\end{description}

\subsection*{Užitečné balíčky pro \LaTeX}

Studenti při sazbě textu často řeší stejné problémy. Některé z nich lze vyřešit následujícími balíčky pro \LaTeX:

\begin{itemize}
  \item \verb|amsmath| -- rozšířené možnosti sazby rovnic,
  \item \verb|float, afterpage, placeins| -- úprava umístění obrázků,
  \item \verb|fancyvrb, alltt| -- úpravy vlastností prostředí Verbatim, 
  \item \verb|makecell| -- rozšíření možností tabulek,
  \item \verb|pdflscape, rotating| -- natočení stránky o 90 stupňů (pro obrázek či tabulku),
  \item \verb|hyphenat| -- úpravy dělení slov,
  \item \verb|picture, epic, eepic| -- přímé kreslení obrázků.
\end{itemize}

Některé balíčky jsou využity přímo v šabloně (v dolní části souboru fitthesis.cls). Nahlédnutí do jejich dokumentace může být rovněž užitečné.

Sloupec tabulky zarovnaný vlevo s pevnou šířkou je v šabloně definovaný \uv{L} (používá se jako \uv{p}).

 % viz. prilohy.tex / see prilohy.tex
\end{document}
